% Pregunta de Investigación
\section{Pregunta de Investigación}

\begin{frame}
\frametitle{Pregunta central}
\begin{center}
\large
\textbf{¿Qué factores sociodemográficos observables se asocian con la probabilidad de reportar consumo de marihuana en los últimos doce meses en Colombia, y cuánto mejora la predicción al pasar de un benchmark econométrico a modelos de \textit{machine learning}?}
\end{center}
\vspace{0.4cm}
\normalsize
Componentes:
\begin{enumerate}
    \item \textbf{Descriptivo-causal}: signo, magnitud y estabilidad de los predictores.
    \item \textbf{Predictivo}: ganancia del ML sobre el lineal manteniendo el mismo \textit{feature set}.
\end{enumerate}
\end{frame}

\begin{frame}
\frametitle{Variable objetivo y unidad de análisis}
\textbf{Variable dependiente:} \(\texttt{propension\_12m}\in\{0,1\}\) — reportó consumo de marihuana en los últimos 12 meses (base limpia, derivada de K\_03 + señales auxiliares).

\vspace{0.3cm}
\textbf{Unidad de análisis:} individuo encuestado (ENCSPA).

\vspace{0.3cm}
\textbf{Tipo de problema:} clasificación binaria.

\vspace{0.4cm}
\textit{Nota crítica}: la construcción de \texttt{propension\_12m} difiere de \texttt{K\_03} \textit{raw} en 898 casos. Se documenta como limitación y se atiende con un análisis de robustez (Y estricta = K\_03=1).
\end{frame}

\begin{frame}
\frametitle{Objetivos}
\textbf{General:} construir un pipeline reproducible que compare benchmark econométrico y \textit{machine learning} para la propensión al consumo, con interpretabilidad y reflexión ética explícitas.

\vspace{0.4cm}
\textbf{Específicos:}
\begin{enumerate}
    \item Integrar y limpiar la base de la ENCSPA con módulos sociodemográficos.
    \item Estimar benchmark MCO/Probit y modelos ML comparables.
    \item Validar con CV-10 y tunear hiperparámetros del modelo final.
    \item Interpretar con valores SHAP y verificar robustez.
    \item Discutir sesgos, ética y usos legítimos del modelo.
\end{enumerate}
\end{frame}
